\section{ROAR: reluzamiento 3D}

\subsection{Motivación: control de grano fino}

Además de la generación del contenido 3D en sí, el control de grano fino sobre la salida es también crucial; \textbf{iluminación} es una de las dimensiones de control más importantes. El problema que resuelve el método ROAR es: dado un conjunto de imágenes con poses y una condición de iluminación objetivo, generar el modelo 3D bajo dicha iluminación.

\subsection{Diseño del método}

ROAR entrena un \textbf{modelo de reluzamiento multivista}, cuya arquitectura es similar a los modelos multivista en CAT3D y B3D, pero que añade la condición de iluminación como entrada.

En concreto:
\begin{itemize}
    \item Entrada: un conjunto de imágenes, sus poses correspondientes + condición de iluminación objetivo
    \item Salida: imágenes reluzadas para cada vista bajo la iluminación objetivo
    \item Objetivo: para $N$ vistas y $M$ condiciones de iluminación, generar $N \times M$ imágenes reluzadas
    \item Estas imágenes se utilizan posteriormente para entrenar un \textbf{NeRF con condición de iluminación}
\end{itemize}

La idea central es similar a un escenario de luz virtual (light stage): al simular la apariencia del objeto bajo diferentes condiciones de iluminación, se entrena al NeRF para generalizar a cualquier iluminación objetivo.

\subsection{Representación de la iluminación: mapa de entorno + incrustación de reflexión especular}

La codificación de la iluminación es un problema no trivial. ROAR adopta una representación dual de la iluminación:

\begin{importantbox}{Codificación dual de la iluminación en ROAR}
\textbf{1. Incrustación de iluminación global}: se convierte el mapa de entorno (environment map) en un vector de incrustación de iluminación mediante un codificador Transformer.

\textbf{2. Incrustación de reflexión especular}: dado que el codificador global no es lo suficientemente potente para capturar completamente los detalles del brillo especular, ROAR introduce adicionalmente una incrustación de reflexión especular (specular embedding), realizando una muestra local del mapa de entorno desde la dirección de reflexión especular. Para manejar diferentes grados de rugosidad, se utiliza un mapa de entorno preborroso, de modo que solo sea necesario muestrear una única posición para obtener información de reflexión especular.
\end{importantbox}

\subsection{Resultados experimentales}

Los experimentos demuestran que ROAR es capaz de recuperar fielmente los efectos del brillo especular bajo diferentes condiciones de iluminación objetivo, mientras que los métodos competitivos omiten el brillo especular o lo compensan en exceso.

\subsection{Resumen de la sección}

ROAR extiende el paradigma de «generación 2D + destilación 3D» a la dimensión del control de iluminación, generando datos ricos de iluminación-vista mediante el entrenamiento de un modelo de reluzamiento multivista y destilándolos en un NeRF con condición de iluminación, logrando así el reluzamiento arbitrario de objetos 3D.


